\section{Discussion}
\label{sec:discussion}

\subsection{Principal findings}

Adherence to the national inpatient hyperglycemia recommendations improved
after their introduction, and higher adherence was consistently associated with
lower mortality and reduced hyperglycemic complications. Adherence was not
uniform: near-universal implementation was observed for several admission and
medication-management actions, whereas the largest gaps were concentrated in
insulin management. This pattern separates recommendations satisfied by a single
discrete action, such as ordering a laboratory test or discontinuing a
medication on admission, from those requiring repeated insulin decisions
across the hospitalization. The improvement was significant but modest, its $P$
value reflecting the number of admissions rather than the size of the effect,
and it was accompanied by substantial inter-hospital heterogeneity that the data
do not explain.

The clearest outcome signal was for 30-day mortality. The propensity-weighted
estimate is primary, being the only one whose exposure window---a fixed
early-admission period---is guaranteed to precede the outcome. Two descriptive
estimates agree in direction: whole-admission adherence adjusted for baseline
characteristics, and a graded fall in mortality across the adherence range. A
graded association is harder to explain by a simple difference in case mix,
though it shares the whole-admission timing concern noted below. Because the
three use different exposure definitions, they corroborate direction rather
than magnitude.

The glycemic findings show a favorable risk profile. Greater adherence bought a
large reduction in hyperglycemic burden---roughly ninefold larger on the
log-odds scale than the accompanying rise in hypoglycemia---at the cost of a
small increase in mild hypoglycemia and no detectable increase in severe
hypoglycemia, which is the balance the recommendations are intended to strike.
The same asymmetry appears at component level, where insulin following
hyperglycemia tracked lower severe hyperglycemia alongside more hypoglycemia.
Since nonsignificance is not evidence of absence, the severe-hypoglycemia
result shows no signal of harm at this sample size rather than establishing
none exists.

Adherence features ranked among the strongest retained predictors alongside
established clinical variables---the practical requirement for treating them
as monitorable quality targets. These findings have no direct precedent: the
position paper was issued in 2023, and no prior study has measured adherence
to it at this scale.

\subsection{Implications}

The most actionable pattern in these results is a convergence. The
recommendations showing the largest implementation gaps---insulin following
hyperglycemia, and the frequency and dose components of routine basal
insulin---are also among those most consistently associated with subsequent
outcomes after balancing measured baseline characteristics. The components
where practice departs most from the guideline are therefore not peripheral
ones, and this identifies them as natural targets for decision support: they
are repeated insulin decisions distributed across a hospitalization,
precisely the kind of task where a prompt at the point of care can plausibly
help, and unlike the admission recommendations they are not already at
ceiling.

This should not be read as an attribution of fault. A low adherence score
records that the action expected by the guideline is absent from the record,
not that the decision was wrong. Ward-level judgment---a patient not eating, an
imminent procedure, a deliberate decision to tolerate higher glucose in a frail
patient---is frequently the correct reason for a departure, and none of it is
recoverable from records collected for clinical rather than research purposes.
The measure is a population-level statistic, as is the inter-hospital variation
reported here. What the work offers is a
measurement instrument---an automated, component-level and transferable way to
ask which of a guideline's recommendations were followed and which track
outcomes---rather than a decision-support system.

Nothing restricts the approach to recommendations a guideline already contains.
Because each is declared externally as a pattern rather than implemented in
analysis code, a practice merely suspected to matter costs no more to evaluate
than the patterns scored here. That inverts the question: one could ask which
candidate practices earn a place in a guideline, and---more
uncomfortably---which current recommendations do not track outcomes at all.
Such candidates would be hypotheses for prospective evaluation, and a wide
search would need more conservative multiplicity correction than applied
here.

\subsection{Limitations}

The design is observational, and none of the estimates should be read as causal
effects. Residual confounding remains throughout: weighting balances only
measured covariates, and unmeasured factors---real-time illness severity above
all, and decisions to limit or de-escalate treatment---cannot be excluded.
Direction was not uniform across outcomes, reflecting differences in indication
and applicability.

The mortality analyses divide by exposure definition. The component-level and
propensity-weighted models score adherence over a fixed early-admission window,
so exposure strictly precedes outcome. The adjusted estimate instead averages
adherence over the whole admission and is descriptive: because 30-day mortality
is predominantly an end-of-admission event, and guideline-concordant management
is frequently de-escalated in patients on a deteriorating trajectory, that
average incorporates the period immediately preceding death, where adherence
may reflect the course of care rather than determine it. Dichotomizing
adherence at the median, as the weighting requires, discards dose-response
information and costs power.

Adherence and outcomes were ascertained during the same hospitalization, and
detecting glycemic events depends on glucose being measured; routine glucose
monitoring was therefore excluded from the glycemic models, though residual
detection effects may remain for other monitoring recommendations. The
pre--post comparison is uncontrolled, so secular trends in inpatient care
cannot be separated from the implementation itself, and the non-mortality
models do not account for the competing risk of death, which biases those
endpoints conservatively. The findings come from adult inpatients in Israeli
general hospitals under a single national guideline, so the adherence levels
observed and the components that proved most consequential may differ
elsewhere. Finally, the data are assembled from hospital admissions rather than
from a longitudinal record, so chronic diagnoses managed in the community may
be incompletely captured; ascertaining eligibility and glycemic outcomes from
glucose measurements rather than diagnosis codes limits the consequences, but
the comorbidity score and the proportion of hyperglycemic admissions carrying a
documented diagnosis are both likely understated.
